% ============================================================
% ch27.tex —— 第 27 章 蚂蚁几何学：高斯的绝妙定理
% ============================================================
\chapter{蚂蚁几何学：高斯的绝妙定理}

上一章我们知道了球面弯、平面平。但"弯"这个量，一直是\textbf{从外面}看的——我们站在三维空间里，俯视球面，看到它瘪下去。这一章回答一个更狠的问题：

\begin{kaiwei}
一只蚂蚁，一生只被允许贴着表面爬，永远不能起飞、不能跳、不能从外面看自己的世界一眼。它能知道自己的世界是弯的吗？
\end{kaiwei}

这个问题不是脑筋急转弯。它之所以致命，是因为\textbf{我们全都是那只蚂蚁}：人类被束缚在三维空间的"表面"（确切说，四维时空的"里面"），不能跳到宇宙外面看它弯不弯。蚂蚁如果测不出自己世界的弯曲，我们也测不出——宇宙学就地解散。蚂蚁如果测得出，宇宙学就有方法论。

\section{蚂蚁的工具箱}

先清点蚂蚁能做什么。它不能"看"曲面的整体形状（那是外在视角），它能做的只有\textbf{表面上的操作}：
\begin{itemize}
\item \textbf{量距离}：沿表面爬，数步数（测地距离，第 26 章的"绷紧线"）；
\item \textbf{量角度}：两条路径的夹角（角度是局部量，表面内完全可测）；
\item \textbf{画圆、画三角形}：从定点出发走等测地距离，把终点连起来。
\end{itemize}
三项全出土的"表面测量"。第 26 章的烧脑时刻已经预演了关键实验（球面圆缩水），现在把它扶正成蚂蚁的\textbf{曲率计}：

\section{蚂蚁的曲率计：画圆，量周长}

蚂蚁在表面上取定点 $O$，向四面八方走出测地距离 $\rho$，把所有终点连成一条"圆"。然后量周长 $C$，与平面世界的标准答案 $2\pi\rho$ 对账：

\begin{center}
\begin{tabular}{ccc}
\toprule
世界 & 蚂蚁的发现 & 判决 \\
\midrule
球面（正弯） & $C < 2\pi\rho$（圆圈"缩水"） & 越走圈得越紧 \\
平面（零弯） & $C = 2\pi\rho$（精确成立） & 标准答案 \\
马鞍面（负弯） & $C > 2\pi\rho$（圆圈"虚胖"） & 越走散得越开 \\
\bottomrule
\end{tabular}
\end{center}

\begin{tikzpicture}[scale=1.0]
  \draw[fill=molv!8] (0,0) .. controls (1.2,0.35) and (2.4,-0.35) .. (3.6,0) .. controls (2.4,0.35) and (1.2,-0.35) .. (0,0);
  \node at (1.8,-0.75) {\scriptsize 薯片/马鞍面：负弯，$C > 2\pi\rho$};
  \draw[fill=shenlan!8] (5.2,0) circle (0.75);
  \draw[fill=white] (5.2,0) circle (0.45);
  \node at (5.2,-1.05) {\scriptsize 球面：正弯，$C < 2\pi\rho$};
  \draw[fill=cheng!10] (7.4,-0.25) rectangle (9.0,0.25);
  \node at (8.2,-0.75) {\scriptsize 平面：零弯，$C = 2\pi\rho$};
\end{tikzpicture}

用第 26 章的公式验算球面那一行（半径 $R$ 的球，测地半径 $\rho$）：圆的三维半径是 $R\sin\rho$（$\rho$ 以弧度计），周长 $C = 2\pi R\sin\rho$；期望值 $2\pi\rho\times R$。因 $\sin\rho < \rho$，\textbf{缩水确凿}。马鞍面的验算反向同理（可用一个具体马鞍面模型——如 $z = x^2 - y^2$ 的原点附近——算出周长反而虚胖，计算放烧脑时刻）。

精确定量：缩水/虚胖的幅度与弯曲程度成正比。小半径展开式（泰勒首项，$\sin\rho = \rho - \tfrac{\rho^3}{6} + \cdots$）：
\[
C = 2\pi R\sin\rho = 2\pi R\rho\left(1 - \frac{\rho^2}{6R^2} + \cdots\right) = 2\pi\rho\left(1 - \frac{K\rho^2}{6} + \cdots\right),
\qquad K = \frac{1}{R^2}.
\]
\textbf{缺口正比于 $K$}——量缺口，即量 $K$。这就是蚂蚁的曲率计。

\begin{faming}
\textbf{土名}：世界的弯曲度 $\to$ \textbf{正式名}：\textbf{高斯曲率} $K$。球面 $K = +\tfrac1{R^2} > 0$，平面 $K = 0$，马鞍面 $K < 0$（大小随点变化）。蚂蚁圆周实验是它的\textbf{操作定义}：弯曲是"用表面测量就能读出的数"。
\end{faming}

\section{圆柱骗局：弯，但蚂蚁量不出}

主菜之前，先上一个思想实验当冷盘——它逼出一个深刻的区分。

把一张纸\textbf{卷成圆柱}。纸"弯"了吗？从外面看，当然弯了（外形从平面变成管子）。现在让蚂蚁在圆柱面上做全套实验：画测地圆、量周长、画三角形、量角和——它会得到什么结果？

\textbf{分毫不差的平面答案}：$C = 2\pi\rho$ 精确成立，内角和恰 $180^\circ$。为什么？注意卷纸的关键事实：\textbf{没有任何两点的表面距离被改变}（纸没拉伸、没褶皱，把圆柱剪开摊平，纸面完好如初）。而蚂蚁的全部实验只用到"表面距离"——距离没变，读数就没变。\textbf{蚂蚁判定：我的世界是平的。}它是对的！圆柱的 $K = 0$（下一节的高斯定理保证）。

于是"弯"分裂成两种，这条区分是微分几何的命脉：

\begin{center}
\begin{tabular}{lp{0.36\textwidth}p{0.36\textwidth}}
\toprule
 & 外在的弯 & 内在的弯 \\
\midrule
是什么 & 形状在空间里"摆放"得弯不弯（纸卷没卷） & 表面本身被拉伸、压缩的程度 \\
谁能看 & 只有外面的三维观察者（蚂蚁看不见） & 蚂蚁用表面测量就能读出 \\
例子 & 圆柱、圆锥的外观 & 球面、马鞍面 \\
改变它 & 卷起来就有了，摊开就没了 & 不拉伸摊不平，撕开才摊得平 \\
\bottomrule
\end{tabular}
\end{center}

高斯证明了数学史上最令人拍案的结果之一：\textbf{高斯曲率 $K$ 完全是内在的}。

\begin{dingli}[绝妙定理（Theorema Egregium），1827]
高斯曲率在"保距变形"（不拉伸、不撕裂、只许弯摆）下\textbf{不变}；它可以通过纯表面测量（蚂蚁的实验）完全确定。
\end{dingli}

名字是高斯自己起的，拉丁文 egregium 意为"卓越超群"——一贯谦虚到近乎刻板的高斯，为此定理破了例。它的物理对应是：\textbf{弯曲是内在属性，不是"住在几维空间里的外观"}。这也一锤定音地解释了地图难题：

\begin{jielun}
球面（$K>0$）与平面（$K=0$）内在曲率不同，故\textbf{不存在}保距的摊平——\textbf{任何世界地图必然撒谎}。
\end{jielun}

谎言的具体形状：墨卡托投影为保持"指南针方向"（航海刚需），把高纬度横向大幅拉伸——格陵兰岛（$216$ 万平方公里）被画到与非洲（$3037$ 万平方公里，相差 \textbf{14 倍}）视觉相当；南美洲被画得只有格陵兰的六成，实际是它的八倍。看世界地图先看投影方式——这堂课值一张机票钱。而"橘子皮永远摊不平、一摊就裂"这句厨房格言，就是绝妙定理的餐桌版。

\section{薯片、甜甜圈与内在曲率的日常}

\textbf{负曲率}世界长什么样？拿起一片品客薯片：中间下凹（一个方向弯向下）、两沿上翘（垂直方向弯向上）——两个方向的弯曲\textbf{一正一反}，乘起来 $K < 0$。它的标志性格局：\textbf{没办法平贴桌面}，四角永远翘。这不是工厂手艺问题，是 $K<0$ 的必然：马鞍面摊平必撕裂（与球面同罪，方向相反）。顺便，为什么薯片设计成马鞍形？\textbf{负曲率结构抗弯强度极高}（两个方向的曲率互相牵制），易碎的薄片因此能承受可观的压力——工程里叫"双曲面壳"，发电厂冷却塔就是放大版薯片。

\textbf{甜甜圈（环面）}则是一座曲率的主题公园：朝外的鼓起部分 $K > 0$，朝内圈的凹陷部分 $K < 0$，两者之间必有一圈\textbf{$K = 0$ 的分界线}（从正到负，连续函数过零）——蚂蚁在环面上走一圈，能依次测出三种世界。还有一条全局预算：环面的"总曲率" $\int K\,dA = 0$（正负恰好抵消）——球面则是 $4\pi$（第 26 章角盈对账的正式版）。\textbf{总曲率被表面的"洞数"钉死}，与具体形状无关——这条"高斯-博内定理"直通拓扑学（第 31 章的莫比乌斯带是近亲）。

\begin{dongshou}
\begin{itemize}
  \yisi 圆柱面上的蚂蚁画测地圆，证明 $C = 2\pi\rho$ \textbf{精确}成立。（把圆柱沿母线剪开摊平——保距变形！蚂蚁的一切测量在摊平后照抄平面答案。这同时证明了圆柱的 $K = 0$：一张 $K=0$ 的纸，怎么卷都还是 $K=0$。）
  \yier 拿一个真甜甜圈（或橡皮泥捏的环面），凭手感找出 $K > 0$（外圈鼓起）、$K < 0$（内圈凹槽）、$K = 0$（分界圈）三个区域，并在表面画出 $K=0$ 的那条闭合曲线。
  \yier 用蚂蚁圆周公式反推地球：蚂蚁在地球表面画 $\rho = 100$ km 的测地圆。先算缩水率 $\tfrac{K\rho^2}{6} = \tfrac{\rho^2}{6R^2}$：$\rho/R = 100/6371 \approx 0.0157$（弧度），其平方除以 $6$ 得 $4.1\times10^{-5}$。缺口 $= 2\pi\rho\times4.1\times10^{-5} \approx 628.3\,\text{km}\times4.1\times10^{-5} \approx 26$ 米——一百公里周长的圆"缩水"二十六米，现代大地测量轻松捕捉。\textbf{蚂蚁圆周计在真地球上实测成功}。再缩到操场尺度（$\rho = 50$ m）：缺口只剩约 $3\times10^{-9}$ 米——三个原子的宽度，任何仪器都无从察觉——这就是你在操场上从未发现地球是圆的的原因。
  \yisi 判断下列各物的 $K$ 符号：鸡蛋壳（各点 $>0$）；自行车内胎（外正内负）；拧过的纸带（"可展"，$K=0$，处处可摊平）；伞面（一般点 $K=0$，伞尖例外——尖点是曲率炸弹）。
\end{itemize}
\end{dongshou}

\begin{shaonao}
\textbf{马鞍面的圆周实验（亲手算一次负曲率）。}取马鞍面 $z = x^2 - y^2$ 的原点附近，蚂蚁画测地圆：

(1) 曲面在原点附近，沿 $x$ 方向的剖面是开口向上的抛物线（弯向上，曲率 $+2$，第 25 章顶点公式），沿 $y$ 方向开口向下（曲率 $-2$）。

(2) 蚂蚁沿 $x$ 轴走距离 $s$，水平位移约 $s$（竖直增量是 $s^2$ 量级，小）；但沿"斜方向"走时，高度先升后降，路径被"折叠"得更短——精确计算（用曲面第一基本形式 $E = 1 + 4x^2$, $G = 1 + 4y^2$ 在原点附近展开）给出：测地半径 $\rho$ 的圆，周长
\[
C = 2\pi\rho\left(1 + \frac{\rho^2}{6} + \cdots\right) > 2\pi\rho \qquad (\text{虚胖}) .
\]
(3) 与球面公式对比：$1 - \tfrac{K\rho^2}{6}$ 中的 $K$，球面是 $+\tfrac1{R^2}$，马鞍面原点是 $-4\times(-1) = -4$（两个方向主曲率 $+2$ 与 $-2$ 之积）——\textbf{正曲率缩水、负曲率虚胖、乘积定符号}。

不想算可以只记结论；算过一遍的人，从此"高斯曲率 $=$ 两方向主曲率之积"不再是口号。深水区的完整公式（$K = \tfrac{LN-M^2}{EG-F^2}$）见 do Carmo 第六章。
\end{shaonao}

\begin{chengshi}
圆周公式 $C = 2\pi\rho(1 - K\rho^2/6 + \cdots)$ 是泰勒展开首项的近似（$\rho$ 小时精确，大圆需高阶项）；绝妙定理的证明（$K$ 可由"第一基本形式"及其导数表出，与嵌入方式无关）超出本书——它是 do Carmo 曲面论的高潮，需要一套本书没铺的记号。"圆柱上 $C = 2\pi\rho$ 精确成立"依赖测地圆恰为横截圆的对称性，一般圆柱点上的圆盘更绕，但摊平论证的精神不变。地球习题用了理想球，实际测量还有地形与重力位的修正。
\end{chengshi}

蚂蚁手里有了曲率计，且这曲率计是\textbf{内在的}——不依赖"外面的世界"。最后一章：把平行线搅进浑水，三种几何逐鹿宇宙——谁赢了？答案是：\textbf{实验说了算}。而我们，就是实验本身。
